
\documentclass[11pt]{article}
\usepackage[a4paper,margin=1in]{geometry}
\usepackage{fontspec}
\setmainfont{DejaVu Serif}
\setsansfont{DejaVu Sans}
\setmonofont{DejaVu Sans Mono}
\usepackage{amsmath,amssymb,mathtools}
\usepackage{microtype}
\usepackage[hidelinks]{hyperref}
\usepackage{enumitem}
\usepackage{parskip}
\allowdisplaybreaks
\setlength{\emergencystretch}{3em}
\newcommand{\dd}{\,d}
\newcommand{\eps}{\varepsilon}
\newcommand{\pder}[2]{\frac{\partial #1}{\partial #2}}
\newcommand{\td}[1]{\frac{d #1}{dt}}
\newcommand{\tdd}[1]{\frac{d^2 #1}{dt^2}}
\newcommand{\detm}[1]{\det\!\left(#1\right)}
\title{Dirichlet XXV: Investigations on a Problem of Hydrodynamics}
\author{G. Lejeune Dirichlet; prepared from the Nachlass by R. Dedekind}
\date{English translation, Werke Band II, pp. 263--302}
\begin{document}
\maketitle

\begin{center}
{\Large \textbf{XXV. Investigations on a Problem of Hydrodynamics}}\\[0.5em]
G. Lejeune Dirichlet\\
Prepared from Dirichlet's papers by R. Dedekind\\
From the eighth volume of the Transactions of the Royal Society of Sciences at Göttingen.
\end{center}

\section*{Preface}

A few remarks must be made in advance about the completion and publication of this memoir, which by the author's last will was entrusted to me.  The hydrodynamical problem treated here, whose solution dates from the winter of 1856--1857, was first presented in brief outline at the close of Dirichlet's lectures on partial differential equations in July 1857.  At the same time the principal result of the whole investigation was published in a short notice in the \emph{Nachrichten} of the Royal Society of Sciences.  The complete exposition was delayed, partly because the author wished to investigate the object still further in its details, and partly because he was occupied with other work, until sudden illness and premature death made completion impossible.

Among the papers that referred to this subject and that came into my hands on 21 July 1859 there was first a manuscript executed with such care that it could be given to the printer without the least alteration.  It is only very much to be regretted that even in this fragment the introduction, which was devoted to the discussion of some general properties of the fundamental equations of hydrodynamics, was left unfinished.  Besides this manuscript, which in the following arrangement extends almost to the end of \S3, there was a large number of separate papers with formulas rapidly written down and without text; their meaning, however, was easy to recognize.  For the most part they were repetitions of what was already presented, and only rarely did one of them give a point of departure for a further development.

With the aid of these papers it was not difficult to find the seven first-order integral equations mentioned in the preliminary notice; they are found in \S5 below.  Moreover, many passages pointed to the case treated in \S8, although nowhere was there a discussion of it; I have tried in \S6 to connect it with the other case investigated in \S7, which, because of its simplicity, was also communicated in the preliminary notice.  Furthermore, as may be seen from the notes that have all been added by me, some passages of the manuscript occasioned computations more laborious than difficult.  Since their results may be useful for future work, I have included them in the memoir, and they form \S4.  After deriving them, I used them in some further investigations; the results, insofar as they have succeeded up to now, are communicated in the final paragraph.

I do not conceal from myself that, despite all the care and affection spent on the work, many things could have been carried out more completely and better.  But I did not wish to delay publication still longer, all the less because I may trust that this last work of the great thinker, who was not granted the opportunity to put the master's hand to the exposition himself, will be appreciated even in its imperfect form.

\begin{flushright}
Zurich, 10 November 1859.\\
R. Dedekind.
\end{flushright}

\section*{Introduction}

In establishing the general equations by which the motion of fluid bodies is determined, one may begin from two different points of view.  According to the first way of conceiving the subject, one sets oneself the task of determining, for an arbitrary place \((x,y,z)\) and an arbitrary time \(t\), the state of the moving mass: its density, its pressure, and the three components of its velocity.  These five quantities are then to be determined as functions of the four variables \(x,y,z,t\).  To this point of view correspond the fundamental equations of hydrodynamics that occur in all textbooks and that Euler first set up.  Euler's equations also form the basis of a great memoir that Lagrange published more than twenty years later in the same academic collection, and from which he later formed, with some additions, the section of his \emph{Mécanique analytique} devoted to hydrodynamics.

The most important of these additions stands at the beginning of that section and concerns a treatment essentially different from Euler's.  Lagrange seeks to follow the motion of each element of the fluid; that is, he determines the coordinates \(x,y,z\), the pressure, and the density of this element by means of its initial coordinates \(a,b,c\) and the time \(t\) elapsed since the beginning of the motion.  Remarkably, however, Lagrange makes no use of the equations corresponding to this point of view.  After observing that they are somewhat complicated, he transforms them into Euler's equations and then adds that the latter, because of their greater simplicity, are preferable for the solution of special problems.

Dirichlet's objection is that the greater formal simplicity of Euler's equations is more than counterbalanced by a hidden difficulty.  In Euler's form the coordinates \(x,y,z\) are not independent variables in the proper sense of the word, since the extent in which they are valid is the region occupied at the given instant by the moving mass; that region is itself determined by the whole preceding motion.  We now know how essential a part of a problem is the domain assigned to functions of several variables that are defined by partial differential equations and boundary conditions.  Euler's form is therefore at its best when the fluid mass preserves the same exterior shape during the motion, or when a fixed solid body moves in an infinite fluid.  The problem considered here is of the opposite type: the boundary is free and changes with the motion.

Dirichlet also points out an error, protected by great authority and therefore repeated in textbooks.  Euler had observed that his equations are greatly simplified if, for the whole duration of the motion, both the three velocity components and the three force components are the partial derivatives, with respect to the coordinates, of functions of those coordinates.  Lagrange added the important assertion that this condition on the velocity components holds automatically for all time if it holds initially and if the force components satisfy the same condition at each time.  Dedekind notes that the manuscript breaks off just here.  The correction intended is essentially this: the property belongs not to a fixed absolute region of space, but to the material mass that is being followed.  If initially, in a region occupied by the fluid, \(u\dd x+v\dd y+w\dd z\) is a complete differential, then the corresponding statement holds later in the region that contains the same material elements.

\section*{\S1. Lagrange's equations and Dirichlet's linear ansatz}

We restrict ourselves to the homogeneous case and put the density equal to one.  Lagrange's fundamental equations are, in Dirichlet's notation,
\begin{align}
 \left(\tdd{x}-X\right)\pder{x}{a}
 +\left(\tdd{y}-Y\right)\pder{y}{a}
 +\left(\tdd{z}-Z\right)\pder{z}{a}
 +\pder{p}{a}&=0,\notag\\
 \left(\tdd{x}-X\right)\pder{x}{b}
 +\left(\tdd{y}-Y\right)\pder{y}{b}
 +\left(\tdd{z}-Z\right)\pder{z}{b}
 +\pder{p}{b}&=0,\\
 \left(\tdd{x}-X\right)\pder{x}{c}
 +\left(\tdd{y}-Y\right)\pder{y}{c}
 +\left(\tdd{z}-Z\right)\pder{z}{c}
 +\pder{p}{c}&=0,\notag\\
 \det\left(\frac{\partial(x,y,z)}{\partial(a,b,c)}\right)&=1.\notag
\end{align}
Here \(a,b,c\) are the initial coordinates of an arbitrary element, \(x,y,z\) its coordinates at the time \(t\), \(p\) the pressure exerted on it, and \(X,Y,Z\) the components of the accelerating force.  The last equation expresses the incompressibility of the fluid.

The force to be considered is the attraction of the whole mass on itself, with elementary attraction inversely proportional to the square of the distance.  If \(V\) is the potential of the fluid at the interior point \((x,y,z)\), and if \(\eps\) is the constant measuring the attraction between two unit masses at unit distance, then
\[
 X=\eps\pder{V}{x},\qquad Y=\eps\pder{V}{y},\qquad Z=\eps\pder{V}{z}.
\]
After substitution, the first three equations of (1) become
\begin{align}
 \tdd{x}\pder{x}{a}+\tdd{y}\pder{y}{a}+\tdd{z}\pder{z}{a}
 -\eps\pder{V}{a}+\pder{p}{a}&=0,\notag\\
 \tdd{x}\pder{x}{b}+\tdd{y}\pder{y}{b}+\tdd{z}\pder{z}{b}
 -\eps\pder{V}{b}+\pder{p}{b}&=0,\\
 \tdd{x}\pder{x}{c}+\tdd{y}\pder{y}{c}+\tdd{z}\pder{z}{c}
 -\eps\pder{V}{c}+\pder{p}{c}&=0.\notag
\end{align}

The investigation is confined to the assumption that \(x,y,z\), as functions of \(a,b,c,t\), contain the three material coordinates only linearly.  Dirichlet uses ``linear'' in the strict homogeneous sense: there is no term independent of \(a,b,c\).  Thus
\begin{equation}
\begin{aligned}
 x&=la+mb+nc,\\
 y&=l'a+m'b+n'c,\\
 z&=l''a+m''b+n''c,
\end{aligned}
\end{equation}
where the nine coefficients depend only on \(t\).  Incompressibility gives
\begin{equation}
 \det\begin{pmatrix}l&m&n\\ l'&m'&n'\\ l''&m''&n''\end{pmatrix}=1.
\end{equation}
For \(t=0\), \(x,y,z\) coincide with \(a,b,c\); hence \(l=m'=n''=1\), while the other six coefficients vanish.

Differentiating (3) with respect to time gives the velocity components
\begin{equation}
\begin{aligned}
 u=\td{x}&=\dot l a+\dot m b+\dot n c,\\
 v=\td{y}&=\dot l'a+\dot m'b+\dot n'c,\\
 w=\td{z}&=\dot l''a+\dot m''b+\dot n''c.
\end{aligned}
\end{equation}
Their initial values are not arbitrary; differentiating the determinant condition and putting \(t=0\) gives
\[
 \dot l_0+\dot m'_0+\dot n''_0=0.
\]

Dirichlet then assumes that the fluid initially fills an ellipsoid whose centre is the origin and whose axes coincide with the coordinate axes:
\begin{equation}
 \frac{a^2}{A^2}+\frac{b^2}{B^2}+\frac{c^2}{C^2}=1.
\end{equation}
Solving (3) for \(a,b,c\) gives
\begin{equation}
\begin{aligned}
 a&=\lambda x+\lambda' y+\lambda'' z,\\
 b&=\mu x+\mu' y+\mu'' z,\\
 c&=\nu x+\nu' y+\nu'' z,
\end{aligned}
\end{equation}
where, since the determinant is one, \(\lambda,\mu,\nu,\ldots\) are just the cofactors of the matrix in (3).  Substitution in (5) gives the equation of the boundary at time \(t\):
\begin{equation}
 \frac{(\lambda x+\lambda' y+\lambda''z)^2}{A^2}
 +\frac{(\mu x+\mu' y+\mu''z)^2}{B^2}
 +\frac{(\nu x+\nu' y+\nu''z)^2}{C^2}=1.
\end{equation}
Thus an initially ellipsoidal boundary remains an ellipsoid, concentric with the original one.  More is true: particles that initially form a concentric, similar, and similarly situated ellipsoid retain the analogous relation to the instantaneous boundary.

It remains to show that (3) can be made to satisfy the equations (2).  The potential of a homogeneous ellipsoid at an interior point is a constant plus a quadratic function of the coordinates.  Therefore, after expressing the current ellipsoid in arbitrary axes and returning to the material coordinates, Dirichlet writes
\begin{equation}
 V=H-La^2-Mb^2-Nc^2-2L'bc-2M'ca-2N'ab,
\end{equation}
where \(H,L,M,N,L',M',N'\) are functions of time.  Since the pressure is to be constant on the boundary, its dependence on \(a,b,c\) must be through the same quadratic form that defines the boundary.  With \(P\) denoting the boundary pressure, one has
\begin{equation}
 p=P+\sigma\left(1-\frac{a^2}{A^2}-\frac{b^2}{B^2}-\frac{c^2}{C^2}\right),
\end{equation}
where \(\sigma\) depends only on time.

Substitution gives ten equations for the ten unknown functions \(l,m,n,l',m',n',l'',m'',n'',\sigma\):
\begin{align}
 l\ddot l+l'\ddot l'+l''\ddot l''&=-2\eps L+\frac{2\sigma}{A^2},&
 m\ddot m+m'\ddot m'+m''\ddot m''&=-2\eps M+\frac{2\sigma}{B^2},\notag\\
 n\ddot n+n'\ddot n'+n''\ddot n''&=-2\eps N+\frac{2\sigma}{C^2},&
 m\ddot n+m'\ddot n'+m''\ddot n''&=-2\eps L',\notag\\
 n\ddot m+n'\ddot m'+n''\ddot m''&=-2\eps L',&
 n\ddot l+n'\ddot l'+n''\ddot l''&=-2\eps M',\\
 l\ddot n+l'\ddot n'+l''\ddot n''&=-2\eps M',&
 l\ddot m+l'\ddot m'+l''\ddot m''&=-2\eps N',\notag\\
 m\ddot l+m'\ddot l'+m''\ddot l''&=-2\eps N',&
 \det\begin{pmatrix}l&m&n\\ l'&m'&n'\\ l''&m''&n''\end{pmatrix}&=1.\notag
\end{align}
Dirichlet observes that one could eliminate \(\sigma\), but he retains the symmetric form.

\section*{\S2. Determination of the accelerations and sign of pressure}

Although the system just displayed satisfies the boundary and field equations and has the same number of equations as unknowns, it is still necessary to show that it actually determines a motion.  Given at some instant finite known values of the nine coefficients and their first derivatives, the equations must determine the second derivatives.  Solving the three equations containing \(\ddot l,\ddot m,\ddot n\), and doing the same for the two other rows, gives each second derivative in the form \(e\sigma+f\), where \(e\) and \(f\) are finite functions of the known data.  Substitution in the differentiated determinant equation gives an equation \(e'\sigma+f'=0\).  The coefficient \(e'\) is a sum of squares of quantities which cannot all vanish because of the determinant condition.  Hence \(\sigma\), and then all accelerations, are determined.

The pressure condition is a further physical restriction.  Formula (9) shows that the factor in parentheses takes all values between zero and one inside the fluid.  If \(\sigma\) is ever negative, the exterior pressure \(P\) must be at least \(|\sigma|\), otherwise the pressure inside would become negative.  Only when \(\sigma\) is never negative may the motion occur in empty space without external pressure.

Dirichlet also determines the initial value of \(\sigma\).  At \(t=0\), adding the first three equations in (10) and using the potential of the initial ellipsoid gives a formula in which \(\sigma_0\) is expressed as a sum of squared differences of the initial velocity coefficients, together with the gravitational contribution.  If the velocity matrix is initially symmetric, i.e. if the entries symmetric about the diagonal are equal, then \(\sigma_0\) is positive.  The proof that it remains positive in that case is postponed to \S5.

\section*{\S3. Decomposition of a linear instantaneous motion}

To obtain a simple intuition for the motion considered in \S1, Dirichlet decomposes the instantaneous linear velocity field into two simpler motions.  After substituting the inverse transformation (6) in (4), the velocity components have the form
\begin{equation}
 u=gx+hy+kz,
 \quad v=g'x+h'y+k'z,
 \quad w=g''x+h''y+k''z.
\end{equation}
The incompressibility condition gives
\[
 g+h'+k''=0.
\]

Every linear velocity field can be decomposed into a strain and a rotation.  For the strain, after a suitable change to axes \(\xi,\eta,\zeta\), the components take the diagonal form
\[
 p=a\xi,
 \qquad q=b\eta,
 \qquad r=c\zeta,
\]
with \(a+b+c=0\) in the incompressible case.  For a pure rotation around an axis through the origin the velocity is
\begin{equation}
 u_1=\gamma y-\beta z,
 \qquad v_1=\alpha z-\gamma x,
 \qquad w_1=\beta x-\alpha y.
\end{equation}
Conversely, the sum of these two fields is any prescribed linear velocity field, and the decomposition is unique.  Dirichlet justifies this by reducing the quadratic form associated with the symmetric part of the linear field to principal axes.  The skew-symmetric part is the rotation.  The restriction \(a+b+c=0\) is exactly the preservation of volume.

\section*{\S4. Auxiliary formulas}

Dirichlet next records the formulas that had only been indicated before.  Let an ellipsoid, referred to arbitrary rectangular axes, be written
\begin{equation}
 Sx^2+S'y^2+S''z^2+2Tyz+2T'zx+2T''xy<1.
\end{equation}
Let \(F\) be the quadratic form on the left and let \(F^*\) denote its adjoint form.  If
\[
 \Delta(s)=\sqrt{\det\begin{pmatrix}
 1+sS&sT''&sT'\\
 sT''&1+sS'&sT\\
 sT'&sT&1+sS''
 \end{pmatrix}},
\]
then the potential for an interior point is obtained by a single integral whose integrand is rational in \(s\) and divided by \(\Delta(s)\).  This is the same result one obtains by first passing to the principal axes and then returning to the original axes; the method of the discontinuous factor gives it more directly.

For the present problem the coefficients of this quadratic form depend on the combinations
\begin{gather*}
 P=l^2+l'^2+l''^2,
 \quad Q=m^2+m'^2+m''^2,
 \quad R=n^2+n'^2+n''^2,\\
 P'=mn+m'n'+m''n'',
 \quad Q'=nl+n'l'+n''l'',
 \quad R'=lm+l'm'+l''m''.
\end{gather*}
These six quantities satisfy the determinant relation
\begin{equation}
 PQR-PP'^2-QQ'^2-RR'^2+2P'Q'R'=1.
\end{equation}
The coefficients \(L,M,N,L',M',N'\) of the potential are then given by definite integrals in \(s\) involving \(P,Q,R,P',Q',R'\).  The source writes these formulas explicitly; the important structural point is that the motion enters the potential only through the six quadratic combinations just displayed, and not separately through the nine entries of the transformation matrix.

The same section writes the coefficients of the velocity field (11) in terms of the transformation matrix.  In modern matrix notation, if
\[
 M(t)=\begin{pmatrix}l&m&n\\l'&m'&n'\\l''&m''&n''\end{pmatrix},
\]
then
\[
 \begin{pmatrix}u\\v\\w\end{pmatrix}=\dot M\begin{pmatrix}a\\b\\c\end{pmatrix}
 =\dot M M^{-1}\begin{pmatrix}x\\y\\z\end{pmatrix}.
\]
Thus the matrix of coefficients \((g,h,k;g',h',k';g'',h'',k'')\) is \(\dot M M^{-1}\).  Since \(\det M=1\), its trace is zero.  The instantaneous rotations about the coordinate axes are
\[
 \alpha=\frac12\left(\pder{w}{y}-\pder{v}{z}\right),\quad
 \beta=\frac12\left(\pder{u}{z}-\pder{w}{x}\right),\quad
 \gamma=\frac12\left(\pder{v}{x}-\pder{u}{y}\right).
\]
These formulas are used both in the determination of \(\sigma\) and in the integral equations.

\section*{\S5. Seven first integrals}

Seven first-order integrals are now set up.  Three of them are obtained directly from (10) by subtracting pairs of equations with the same right-hand side.  In component form they express the constancy of the vorticity combination carried by the material coordinates.  In modern notation the skew-symmetric matrix
\begin{equation}
 \mathcal L=M D\dot M^T-\dot M D M^T,
 \qquad D=\operatorname{diag}(A^2,B^2,C^2),
\end{equation}
has constant entries.  Dirichlet gives the three scalar component equations, and then rewrites them in terms of the spatial velocity derivatives
\[
 \pder{w}{y}-\pder{v}{z},\qquad
 \pder{u}{z}-\pder{w}{x},\qquad
 \pder{v}{x}-\pder{u}{y}.
\]
It follows immediately that the axis of instantaneous rotation is always formed by the same elements of the fluid mass.  If the three displayed quantities vanish at one time, so that there is no rotation, then they vanish throughout the motion.

Since the acting forces come only from the mutual attraction of the elements of the fluid, the principle of areas gives three further integrals:
\begin{equation}
 \int(yw-zv)\dd\tau=\text{const.},\quad
 \int(zu-xw)\dd\tau=\text{const.},\quad
 \int(xv-yu)\dd\tau=\text{const.}.
\end{equation}
To compute them, one returns to the original ellipsoid and uses
\[
 \int a^2\dd\tau=\frac{M A^2}{5},\quad
 \int b^2\dd\tau=\frac{M B^2}{5},\quad
 \int c^2\dd\tau=\frac{M C^2}{5},
\]
while the mixed integrals \(\int ab\dd\tau\), \(\int bc\dd\tau\), \(\int ca\dd\tau\) vanish.  This gives Dirichlet's three explicit component formulas involving \(l,m,n\) and their derivatives.

The seventh integral is the principle of living force.  In the present notation it is
\begin{equation}
 \int\left[\left(\td{x}\right)^2+
 \left(\td{y}\right)^2+\left(\td{z}\right)^2\right]\dd\tau
 =\text{constant}+2\int V\dd\tau.
\end{equation}
The left-hand side becomes a quadratic form in \(\dot l,\dot m,\ldots,\dot n''\); the self-potential integral is evaluated from the ellipsoidal formula.  If the current ellipsoid is referred to its principal axes and has semiaxes \(\mathfrak a,\mathfrak b,\mathfrak c\), then Dirichlet obtains the same integral by computing
\[
 \int V\dd\tau
\]
directly over that ellipsoid.  The result agrees with the expression obtained from the formulas of \S4.  The calculation also shows that the symmetry assumptions made later are compatible with the full differential equations.

\section*{\S6. The rotationally symmetric case}

Because the differential equations (10) are highly complicated, a complete solution can be expected only under special initial assumptions.  Dirichlet first assumes complete symmetry, both of shape and of initial velocity, about a fixed axis.  Thus the initial ellipsoid is a spheroid of revolution.  Taking the symmetry axis to be the \(c\)-axis gives \(B=A\).  The initial velocity must be invariant under every rotation of the \((a,b)\)-plane.  This leads to the form
\begin{equation}
 u=ga+hb,
 \qquad v=-ha+gb,
 \qquad w=-2gc.
\end{equation}
The two parts of the instantaneous motion are therefore a rotational part about the symmetry axis and a strain that expands or contracts in the meridian planes.

The same symmetry persists for all time.  The linear transformation therefore has the form
\begin{equation}
 x=la+mb,
 \qquad y=-ma+lb,
 \qquad z=n''c,
\end{equation}
with incompressibility
\begin{equation}
 (l^2+m^2)n''=1.
\end{equation}
The boundary remains a spheroid of revolution.  The velocity components, when written in spatial coordinates, again show a meridian strain together with a rotation about the axis.  The angular velocity about the \(z\)-axis is proportional to
\[
 l\dot m-m\dot l.
\]

Substitution in the fundamental equations reduces the nine equations to three independent equations for \(l,m,n''\) and \(\sigma\), together with (20).  One first integral is
\begin{equation}
 l\dot m-m\dot l=\frac{\gamma}{n''},
\end{equation}
where \(\gamma\) is constant.  Hence the angular velocity is proportional to the length of the axis of rotation.  The equation (20) allows \(l,m\) to be eliminated.  Introducing the ratio
\begin{equation}
 \alpha=\frac{n''}{D},
\end{equation}
of the axis of rotation to the radius \(D\) of the sphere of equal volume, and putting
\[
 f(\alpha)=\int_0^\infty\frac{\dd s}{(\alpha^2+s)\sqrt{(1+s)^2(\alpha^2+s)}},
\]
Dirichlet transforms the equations into a one-dimensional system for \(\alpha(t)\).  The constants of integration are those fixed by the initial shape and initial velocity.

For later discussion he records the known properties of \(f\).  Explicit evaluation gives different elementary forms according as \(\alpha<1\) or \(\alpha>1\); both give \(f(1)=2\).  As \(\alpha\) tends to zero or infinity, \(f(\alpha)\) tends to zero.  The function has a single maximum at \(\alpha=1\).  Also, the function \(\alpha^2 f'(\alpha)\) has a numerical maximum approximately \(0.2246\), a number already known from the theory of uniformly rotating fluid masses.

\section*{\S7. Motion without rotation}

Dirichlet next considers the special case in which initially, and therefore always, no rotation occurs: \(\gamma=0\).  Suppose first that the fluid starts from rest.  Then the equations reduce to
\begin{equation}
 \left(1+\frac{1}{2\alpha^3}\right)\left(\frac{\dd\alpha}{\dd t}\right)^2
 =2\eps\bigl(f(\alpha)-f(\alpha_0)\bigr),
\end{equation}
up to Dirichlet's normalizing constants.  If \(\alpha_0=1\), the initial shape is a sphere and the solution remains the sphere.  If \(\alpha_0<1\), so that the initial spheroid is oblate, there is a second root \(\alpha_1>1\) of
\[
 f(\alpha_1)=f(\alpha_0).
\]
The quantity \(\alpha\) runs periodically through the interval from \(\alpha_0\) to \(\alpha_1\) and back.  The motion consists of isochronous oscillations in which the fluid passes through the spherical shape and alternately assumes the form of a prolate and an oblate ellipsoid.

If the spheroid does not start from rest, the same character persists as long as the initial velocity is below a certain bound depending on the initial form.  If that bound is exceeded, \(\dot\alpha\) does not change sign after a finite time.  With positive initial \(\dot\alpha\), the spheroid lengthens without bound; with negative initial \(\dot\alpha\), it flattens without bound.  In all of these nonrotating cases, however, the quantity \(\sigma\) never becomes negative.  Therefore these motions are physically possible without an imposed exterior pressure.

\section*{\S8. Motion with rotation}

Assume now that \(\gamma\ne0\), so that rotation is present throughout the motion.  From the properties of \(f\) there is a unique \(\delta\) with \(0<\delta<1\) satisfying the equation corresponding to the balance of the rotational term and the gravitational term.  Dirichlet introduces
\[
 \psi(\alpha)=\frac{\gamma^2}{\alpha^2}+\eps f(\alpha),
\]
again up to the constants used in the source.  This function decreases from \(\alpha=0\) to \(\delta\), reaches a negative minimum, and then increases without bound.  The equation of motion may be written in the form
\begin{equation}
 \left(1+\frac{1}{2\alpha^3}\right)
 \left(\frac{\dd\alpha}{\dd t}\right)^2
 =\psi(\alpha_0)+K-\psi(\alpha),
\end{equation}
with \(K\ge0\) determined by the initial velocity.

It follows first that \(\alpha\) always remains below an assignable finite limit: even the smallest initial rotation prevents unlimited elongation of the spheroid.  Three cases remain.

\begin{enumerate}[label=\arabic*)]
\item The constant has the minimum value.  Then \(K=0\), \(\alpha_0=\delta\), and the spheroid keeps its form while rotating uniformly about its short axis.  This is the case first treated by Maclaurin.  The square of the angular velocity must not exceed the known numerical limit \(0.2246\ldots\).  For each value below this bound there are two corresponding spheroids, which coincide when the limiting value is reached.

\item The constant is between the minimum and zero.  The equation \(\psi(\alpha)=\psi(\alpha_0)+K\) has two positive roots \(\alpha'<\alpha''\), with \(0<\alpha'<\delta\).  The variable \(\alpha\) oscillates periodically between these two values.  At the minimum value the rotation is too small, and at the maximum value too large, for the fluid to maintain the instantaneous form.

\item The constant is positive.  Then the same equation has only one positive root, and either from the beginning or after a finite time the spheroid begins to flatten more and more without bound.  In the rotating cases, if the rotation at the instant of greatest elongation is sufficiently large, the physical possibility of the motion may require an exterior pressure.
\end{enumerate}

\section*{\S9. Further ellipsoidal motions and Jacobi's theorem}

The preceding cases have the special property that the three quantities denoted in \S4 by \(P',Q',R'\) remain zero throughout the motion.  Dirichlet and Dedekind therefore ask whether other motions with the same property are possible.  A careful analysis gives two further classes compatible with the fundamental equations.  One of them is still expressible by a diagonal transformation of the axes but is not reducible to quadratures merely by using the integral of living force.

The second class yields Jacobi's beautiful result.  A triaxial ellipsoid whose axes \(A,B,C\) satisfy a certain definite integral relation can rotate uniformly about the smallest axis.  In suitable coordinates the motion is
\begin{equation}
 x=a\cos kt+b\sin kt,
 \qquad y=-a\sin kt+b\cos kt,
 \qquad z=c,
\end{equation}
with the appropriate scaling of axes and with \(k\) fixed by the same integral condition.  This is the Jacobi ellipsoid: a homogeneous self-gravitating fluid ellipsoid can rotate like a rigid body about its centre only when the axis of rotation is fixed and coincides with a principal axis, the Maclaurin and Jacobi cases being the corresponding possibilities.

The geometrical meaning of the equations \(P'=Q'=R'=0\) is that the material elements originally lying on the three coordinate axes, that is, on the principal axes of the ellipsoid, continue throughout the motion to fill three mutually perpendicular straight lines.  Since linear transformations preserve conjugate diameters, the actual meaning is that the three principal axes of the ellipsoid are always formed by the same material elements of the fluid mass.

A related hypothesis is that the directions of the three principal axes remain fixed.  In the notation of \S4 this is expressed by \(T=T'=T''=0\).  It holds in the rotationally symmetric case already treated and in one of the additional cases just mentioned.  It also gives a third case, a remarkable counterpart to Jacobi's theorem.  Every triaxial ellipsoid that satisfies Jacobi's condition can preserve its external shape and position while an internal motion of its elements takes place:
\begin{equation}
 x=a\cos kt+b\frac{A}{B}\sin kt,
 \qquad
 y=-a\frac{B}{A}\sin kt+b\cos kt,
 \qquad
 z=c.
\end{equation}
Each particle describes the ellipse
\[
 \frac{x^2}{A^2}+\frac{y^2}{B^2}=\frac{a^2}{A^2}+\frac{b^2}{B^2},\qquad z=c,
\]
and it does so exactly as if it were isolated and attracted toward the centre of its orbit by a force proportional to distance, whose measure for unit distance is \(k^2\).

\end{document}
